\section{模組說明}
Market 模組提供從市場結構發現到單只 ETF 研究、比較和保存的完整路徑。用戶先在 \texttt{Market Intelligence} 建立市場方向，再進入 \texttt{Sector Detail} 比較板塊內 ETF，隨後在 \texttt{ETF Detail} 深入研究，並將結果保存到 \texttt{Watchlist} 或 \texttt{My Portfolio}。
\begin{quote}\small 資料邊界：頁面展示延遲行情、緩存快照或 EOD（end-of-day）資料時，必須保留相應標簽。Market 模組用於研究發現，不構成交易建議。\end{quote}
\section{User Story /\allowbreak{} Video 索引}
\begingroup
\small
\begin{longtable}{L{0.18\linewidth}L{0.24\linewidth}L{0.28\linewidth}L{0.16\linewidth}}
\toprule
\textbf{Story} & \textbf{視頻主題} & \textbf{核心路徑} & \textbf{Video} \\
\midrule
\hyperlink{us-mkt-01}{US-MKT-01} & Market Intelligence Home & \texttt{Market Intelligence} -\textgreater{} \texttt{Sector Detail} & \href{https://jjpvro70sief.jp.larksuite.com/wiki/WE34wt8oyiO3NOkBA3zjFrxppge}{視頻頁面} \\
\hyperlink{us-mkt-02}{US-MKT-02} & Sector Detail & \texttt{Sector Detail} -\textgreater{} \texttt{ETF Rankings} & \href{https://jjpvro70sief.jp.larksuite.com/wiki/WE34wt8oyiO3NOkBA3zjFrxppge}{視頻頁面} \\
\hyperlink{us-mkt-03}{US-MKT-03} & ETF Detail and Compare ETF & \texttt{ETF Detail} -\textgreater{} \texttt{Compare ETF} & \href{https://jjpvro70sief.jp.larksuite.com/wiki/WE34wt8oyiO3NOkBA3zjFrxppge}{視頻頁面} \\
\hyperlink{us-mkt-04}{US-MKT-04} & Save ETF Research & \texttt{Watchlist} /\allowbreak{} \texttt{Add to My Portfolio} /\allowbreak{} \texttt{Screener} & \href{https://jjpvro70sief.jp.larksuite.com/wiki/WE34wt8oyiO3NOkBA3zjFrxppge}{視頻頁面} \\
\bottomrule
\end{longtable}
\endgroup
\section{統一術語}
\begingroup
\small
\begin{longtable}{L{0.24\linewidth}L{0.68\linewidth}}
\toprule
\textbf{UI 原文} & \textbf{文檔含義} \\
\midrule
\texttt{Market Intelligence} & Market 模組首頁，也是 ETF 研究起點 \\
\texttt{Sector Heatmap} & 按板塊或主題聚合 ETF 表現的熱力圖 \\
\texttt{AI Market Signals} & 將市場變動整理為可閱讀研究信號的卡片區域 \\
\texttt{Market Discovery Rail} & 從板塊、ETF movers、新 ETF、活躍 ETF 等維度繼續發現的橫向入口 \\
\texttt{Sector Detail} & 所選板塊的詳情頁面 \\
\texttt{AI Sector Insights} & 基於 EOD 緩存資料生成的板塊研究觀察 \\
\texttt{ETF Rankings} & 所選板塊內的 ETF 排名列表 \\
\texttt{ETF Detail} & 單只 ETF 的詳情與圖表研究頁面 \\
\texttt{Compare ETF} & 兩只 ETF 的歸一化疊加比較功能 \\
\texttt{Screener} & ETF 搜尋、篩選和排序頁面，包含 \texttt{Recommended}、\texttt{All}、\texttt{Watchlist} 三個範圍 \\
\bottomrule
\end{longtable}
\endgroup
\prdhr
\hypertarget{us-mkt-01}{}
\section{User Story 1 - 從 \texttt{Market Intelligence} 建立研究方向}
\textbf{Story ID:} \texttt{US-MKT-01}
\textbf{用戶故事：} 作為尚未確定具體 ETF 的用戶，我希望先了解當天的市場廣度、板塊變動和主要指數背景，從而選擇下一步值得研究的板塊。
\textbf{Video:} \href{https://jjpvro70sief.jp.larksuite.com/wiki/WE34wt8oyiO3NOkBA3zjFrxppge}{Market 驗收視頻頁面}
\subsection{Walkthrough}
\begin{enumerate}
\item 打開 \texttt{Market Intelligence}。先展示頁面整體結構，說明它是 ETF 研究起點，而不是完整 ETF 目錄或交易推薦頁。
\item 瀏覽 \texttt{Sector Heatmap}。說明每個色塊代表一個板塊或主題籃子，顯示等權變動、ETF 數量和小型趨勢線；同時展示 \texttt{Rising}、\texttt{Flat}、\texttt{Falling} 匯總。
\item 向下瀏覽 \texttt{AI Market Signals}。選擇至少一張信號卡，展示信號日期、市場變動、AI analysis 摘要、相關 ETF，以及存在時的外部來源。
\item 展示 \texttt{Market Indices}，包括 Hang Seng Index、Hang Seng Tech 和 S\&P 500 Index，並讓 delayed data 標簽保持可見。
\item 瀏覽 \texttt{Market Discovery Rail}。在 \texttt{Sectors} 範圍展示 \texttt{Default} 與 \texttt{Hot}，說明它們分別提供穩定入口和當日變動較大的板塊入口。
\item 展示 \texttt{Recommended ETFs} 的 symbol、name、price、daily move 和 \texttt{View all}。
\item 點擊一個板塊（建議使用 \texttt{Crypto}），進入 \texttt{Sector Detail}，完成當前 story 並銜接下一個 story。
\end{enumerate}
\screenshotsTripleCompact{assets/screenshots/mkt-us1-01-sector-heatmap.png}{MKT-US1-01}{\texttt{Market Intelligence} 頂部與完整 \texttt{Sector Heatmap}，包含 \texttt{Rising /\allowbreak{} Flat /\allowbreak{} Falling} 市場廣度匯總。}{assets/screenshots/mkt-us1-02-ai-market-signal.png}{MKT-US1-02}{完整的 \texttt{AI Market Signals} 卡，包含 signal date、AI analysis、sources、related ETFs，以及帶 delayed 標簽的 \texttt{Market Indices}。}{assets/screenshots/mkt-us1-03-discovery-recommended.png}{MKT-US1-03}{\texttt{Market Discovery Rail} 的 \texttt{Default /\allowbreak{} Hot} sector baskets，以及 \texttt{Recommended ETFs} 和 \texttt{View all}。}
\screenshotDescsThree{\textbf{MKT-US1-01} \textemdash{} \texttt{Market Intelligence} 頂部與完整 \texttt{Sector Heatmap}，包含 \texttt{Rising /\allowbreak{} Flat /\allowbreak{} Falling} 市場廣度匯總。}{\textbf{MKT-US1-02} \textemdash{} 完整的 \texttt{AI Market Signals} 卡，包含 signal date、AI analysis、sources、related ETFs，以及帶 delayed 標簽的 \texttt{Market Indices}。}{\textbf{MKT-US1-03} \textemdash{} \texttt{Market Discovery Rail} 的 \texttt{Default /\allowbreak{} Hot} sector baskets，以及 \texttt{Recommended ETFs} 和 \texttt{View all}。}
\subsection{Acceptance Criteria}
\begingroup
\small
\begin{longtable}{L{0.18\linewidth}L{0.30\linewidth}L{0.18\linewidth}L{0.12\linewidth}L{0.10\linewidth}}
\toprule
\textbf{ID} & \textbf{Requirement} & \textbf{Reference} & \textbf{Ownership} & \textbf{Priority} \\
\midrule
MKT-US1-AC01 & \texttt{Market Intelligence} 成功載入，並以 \texttt{Sector Heatmap}、\texttt{AI Market Signals}、\texttt{Market Indices}、\texttt{Market Discovery Rail} 和 \texttt{Recommended ETFs} 組織研究入口。 & US-MKT-01 & Integration & Must \\
MKT-US1-AC02 & \texttt{Sector Heatmap} 的每個可見色塊展示板塊/\allowbreak{}主題名稱、變動、ETF 數量和趨勢線，並通過顏色區分上漲、接近平盤和下跌。 & US-MKT-01 & Integration & Must \\
MKT-US1-AC03 & \texttt{Rising}、\texttt{Flat}、\texttt{Falling} 匯總與當前熱力圖資料對應。 & US-MKT-01 & Integration & Must \\
MKT-US1-AC04 & \texttt{AI Market Signals} 卡展示信號日期、檢測到的變動、AI analysis 摘要和 contributing/\allowbreak{}related ETFs；有外部來源時顯示 source reference。 & US-MKT-01 & Integration & Must \\
MKT-US1-AC05 & \texttt{Market Indices} 展示主要指數及當前可用的變動資料，並清楚顯示 delayed data 邊界。 & US-MKT-01 & Integration & Must \\
MKT-US1-AC06 & \texttt{Market Discovery Rail} 支持從板塊、ETF movers、新 ETF 或活躍 ETF 等範圍繼續研究。 & US-MKT-01 & Integration & Must \\
MKT-US1-AC07 & \texttt{Sectors} 中的 \texttt{Default} 與 \texttt{Hot} 作用域可以切換，且展示內容與當前選擇一致。 & US-MKT-01 & Integration & Must \\
MKT-US1-AC08 & \texttt{Recommended ETFs} 是精選池而非完整 ETF universe；每行展示 symbol、name、price 和 daily move，並提供 \texttt{View all}。 & US-MKT-01 & Integration & Must \\
MKT-US1-AC09 & 點擊一個板塊後，打開與所選板塊匹配的 \texttt{Sector Detail}。 & US-MKT-01 & Integration & Must \\
MKT-US1-AC10 & 頁面沒有把市場信號、板塊表現或精選 ETF 表述為交易建議。 & US-MKT-01 & Integration & Must \\
\bottomrule
\end{longtable}
\endgroup
\prdhr
\hypertarget{us-mkt-02}{}
\section{User Story 2 - 理解板塊走勢並篩選 ETF}
\textbf{Story ID:} \texttt{US-MKT-02}
\textbf{用戶故事：} 作為從市場首頁進入某個板塊的用戶，我希望理解該板塊變動的廣度、流動性與風險背景，並在板塊內比較 ETF，選出下一步研究對象。
\textbf{Video:} \href{https://jjpvro70sief.jp.larksuite.com/wiki/WE34wt8oyiO3NOkBA3zjFrxppge}{Market 驗收視頻頁面}
\subsection{Walkthrough}
\begin{enumerate}
\item 從 \texttt{Sector Heatmap} 或 \texttt{Market Discovery Rail} 打開一個板塊。建議使用資料完整的 \texttt{Crypto}。
\item 展示頂部 sector banner：板塊名稱、代表性 tickers、籃子內 ETF 數量、最新板塊變動，以及 cached sector snapshot 標簽。
\item 瀏覽 \texttt{AI Sector Insights}。依次展示 \texttt{Momentum \& Breadth}、\texttt{Liquidity Pulse}、\texttt{Trend \& Risk}，說明這些內容基於 EOD cached ETF data，不是即時預測。
\item 進入 \texttt{ETF Rankings}。確認列表只包含所選板塊內的 ETF。
\item 依次切換 1 至 2 個排序指標，例如 \texttt{Change}、\texttt{Volume} 或 \texttt{Expense Ratio}；再切換一次 return period，例如 \texttt{1D} 到 \texttt{1M}。
\item 選擇一只 ETF 進入 \texttt{ETF Detail}，完成從板塊發現到產品研究的過渡。
\end{enumerate}
\screenshotsTripleCompact{assets/screenshots/mkt-us2-01-sector-banner.png}{MKT-US2-01}{\texttt{Crypto Sector Detail} banner，包含 representative tickers、ETF count、sector move 和 cached sector snapshot 標簽。}{assets/screenshots/mkt-us2-02-ai-sector-insights.png}{MKT-US2-02}{Consumer sector 的 \texttt{AI Sector Insights}，同時展示 \texttt{Momentum \& Breadth}、\texttt{Liquidity Pulse} 和 EOD data boundary。}{assets/screenshots/mkt-us2-03-etf-rankings.png}{MKT-US2-03}{Crypto \texttt{ETF Rankings}，顯示 \texttt{Change} sort、\texttt{1D} return period 和完整 ETF ranking rows。}
\screenshotDescsThree{\textbf{MKT-US2-01} \textemdash{} \texttt{Crypto Sector Detail} banner，包含 representative tickers、ETF count、sector move 和 cached sector snapshot 標簽。}{\textbf{MKT-US2-02} \textemdash{} Consumer sector 的 \texttt{AI Sector Insights}，同時展示 \texttt{Momentum \& Breadth}、\texttt{Liquidity Pulse} 和 EOD data boundary。}{\textbf{MKT-US2-03} \textemdash{} Crypto \texttt{ETF Rankings}，顯示 \texttt{Change} sort、\texttt{1D} return period 和完整 ETF ranking rows。}
\subsection{Acceptance Criteria}
\begingroup
\small
\begin{longtable}{L{0.18\linewidth}L{0.30\linewidth}L{0.18\linewidth}L{0.12\linewidth}L{0.10\linewidth}}
\toprule
\textbf{ID} & \textbf{Requirement} & \textbf{Reference} & \textbf{Ownership} & \textbf{Priority} \\
\midrule
MKT-US2-AC01 & \texttt{Sector Detail} 展示與用戶所選板塊一致的名稱和資料，不沿用之前打開的板塊內容。 & US-MKT-02 & Integration & Must \\
MKT-US2-AC02 & sector banner 展示代表性 tickers、ETF 數量和最新板塊變動。 & US-MKT-02 & Integration & Must \\
MKT-US2-AC03 & 頁面明確標註 sector snapshot /\allowbreak{} EOD cached data 邊界，不將其描述為即時資料。 & US-MKT-02 & Integration & Must \\
MKT-US2-AC04 & \texttt{AI Sector Insights} 至少覆蓋 \texttt{Momentum \& Breadth}、\texttt{Liquidity Pulse}、\texttt{Trend \& Risk} 三類研究角度。 & US-MKT-02 & Integration & Must \\
MKT-US2-AC05 & 當某項資料不足時，對應 insight 顯示 unavailable/\allowbreak{}data boundary，不虛構結論。 & US-MKT-02 & Integration & Must \\
MKT-US2-AC06 & \texttt{ETF Rankings} 僅顯示當前所選板塊內的 ETF。 & US-MKT-02 & Integration & Must \\
MKT-US2-AC07 & 排序支持當前實現提供的指標，包括 \texttt{Change}、\texttt{Volume}、\texttt{Turnover}、\texttt{Price}、\texttt{Expense Ratio}、\texttt{Beta}、\texttt{Drawdown}、\texttt{Sharpe ratio}、\texttt{Name} 或 \texttt{Symbol}。 & US-MKT-02 & Integration & Must \\
MKT-US2-AC08 & return period 支持在當前實現提供的 \texttt{1D}、\texttt{1W}、\texttt{1M}、\texttt{3M} 之間切換，列表結果與當前周期一致。 & US-MKT-02 & Integration & Must \\
MKT-US2-AC09 & 切換排序指標或周期時，當前板塊範圍保持不變。 & US-MKT-02 & Integration & Must \\
MKT-US2-AC10 & 點擊排名中的 ETF 後，打開與該行 symbol 匹配的 \texttt{ETF Detail}。 & US-MKT-02 & Integration & Must \\
MKT-US2-AC11 & \texttt{AI Sector Insights} 采用研究觀察措辭，不預測收益，也不推薦交易。 & US-MKT-02 & Integration & Must \\
\bottomrule
\end{longtable}
\endgroup
\prdhr
\hypertarget{us-mkt-03}{}
\section{User Story 3 - 研究並比較單只 ETF}
\textbf{Story ID:} \texttt{US-MKT-03}
\textbf{用戶故事：} 作為已選定候選 ETF 的用戶，我希望查看報價、K 線、關鍵統計和新聞，並與另一只 ETF 比較相對走勢。
\textbf{Video:} \href{https://jjpvro70sief.jp.larksuite.com/wiki/WE34wt8oyiO3NOkBA3zjFrxppge}{Market 驗收視頻頁面}
\subsection{Walkthrough}
\begin{enumerate}
\item 從 \texttt{ETF Rankings}、\texttt{Recommended ETFs}、\texttt{All}、\texttt{Watchlist} 或 ETF movers 打開一只資料完整的 ETF。
\item 展示詳情頭部：fund name、symbol、exchange、latest price、currency、daily change 和 market/\allowbreak{}close status。
\item 在 K-line chart 上切換 \texttt{1D}、\texttt{1W}、\texttt{1M}、\texttt{3M} 中至少兩個 time ranges。
\item 在圖表上 scrub，展示所選時間點的 open、high、low、close 和 volume 更新。
\item 打開 \texttt{View} panel，切換一次 chart type，並開啟或關閉至少一條 moving average（\texttt{MA5}、\texttt{MA10} 或 \texttt{MA30}）。
\item 切換 \texttt{Overview} 與 \texttt{News}。在 \texttt{Overview} 核對 key stats；在 \texttt{News} 展示 headline、source 和 timestamp。
\item 打開 \texttt{More}，進入 \texttt{Compare ETF}，選擇第二只 ETF，並查看同一 chart option 下的 normalized overlay。
\end{enumerate}
\screenshotsTripleCompact{assets/screenshots/mkt-us3-01-etf-detail-header.png}{MKT-US3-01}{ChinaAMC Ether ETF 的 \texttt{ETF Detail}，包含 identity、quote、close status、K-line chart、moving averages、Overview 和 key stats。}{assets/screenshots/mkt-us3-02-chart-view-panel.png}{MKT-US3-02}{展開的 \texttt{View} panel，展示 K-line types、當前 \texttt{OHLC bar} 選擇和 \texttt{MA5 /\allowbreak{} MA10 /\allowbreak{} MA30} controls。}{assets/screenshots/mkt-us3-03-compare-etf.png}{MKT-US3-03}{\texttt{ETF Comparison} 完成態，包含 3046.HK 與 3189.HK、normalized price performance、圖例及 factsheet matrix。}
\screenshotDescsThree{\textbf{MKT-US3-01} \textemdash{} ChinaAMC Ether ETF 的 \texttt{ETF Detail}，包含 identity、quote、close status、K-line chart、moving averages、Overview 和 key stats。}{\textbf{MKT-US3-02} \textemdash{} 展開的 \texttt{View} panel，展示 K-line types、當前 \texttt{OHLC bar} 選擇和 \texttt{MA5 /\allowbreak{} MA10 /\allowbreak{} MA30} controls。}{\textbf{MKT-US3-03} \textemdash{} \texttt{ETF Comparison} 完成態，包含 3046.HK 與 3189.HK、normalized price performance、圖例及 factsheet matrix。}
\subsection{Acceptance Criteria}
\begingroup
\small
\begin{longtable}{L{0.18\linewidth}L{0.30\linewidth}L{0.18\linewidth}L{0.12\linewidth}L{0.10\linewidth}}
\toprule
\textbf{ID} & \textbf{Requirement} & \textbf{Reference} & \textbf{Ownership} & \textbf{Priority} \\
\midrule
MKT-US3-AC01 & \texttt{ETF Detail} 展示與所選 ETF 一致的 fund name、symbol、exchange、latest price、currency、daily change 和 market status。 & US-MKT-03 & Integration & Must \\
MKT-US3-AC02 & K-line chart 使用當前 ETF 的價格序列，不顯示先前 ETF 的殘留資料。 & US-MKT-03 & Integration & Must \\
MKT-US3-AC03 & 用戶可在當前實現提供的 \texttt{1D}、\texttt{1W}、\texttt{1M}、\texttt{3M} time ranges 間切換，圖表隨選擇更新。 & US-MKT-03 & Integration & Must \\
MKT-US3-AC04 & scrub 圖表時，open、high、low、close 和 volume 與當前選中資料點同步更新。 & US-MKT-03 & Integration & Must \\
MKT-US3-AC05 & \texttt{View} panel 支持 \texttt{solid candlestick}、\texttt{hollow candlestick}、\texttt{OHLC bar}、\texttt{line} 和 \texttt{area} 中當前已啟用的 chart types。 & US-MKT-03 & Integration & Must \\
MKT-US3-AC06 & \texttt{MA5}、\texttt{MA10}、\texttt{MA30} 開關只改變當前圖表疊加層，不改變 ETF 或 time range。 & US-MKT-03 & Integration & Must \\
MKT-US3-AC07 & \texttt{Overview} 展示當前可用 key stats；不可用值以明確的 unavailable 狀態顯示。 & US-MKT-03 & Integration & Must \\
MKT-US3-AC08 & \texttt{News} 展示與當前 ETF 研究流相關的 headline、source 和 timestamp。 & US-MKT-03 & Integration & Must \\
MKT-US3-AC09 & \texttt{More} 中的 \texttt{Compare ETF} 使用當前 ETF 作為第一只比較對象，並允許選擇第二只 ETF。 & US-MKT-03 & Integration & Must \\
MKT-US3-AC10 & 比較圖使用相同 chart option 載入兩組價格序列，並以 normalized overlay 展示相對表現。 & US-MKT-03 & Integration & Must \\
MKT-US3-AC11 & 退出 \texttt{Compare ETF} 後仍回到原 ETF 的詳情上下文。 & US-MKT-03 & Integration & Must \\
\bottomrule
\end{longtable}
\endgroup
\prdhr
\hypertarget{us-mkt-04}{}
\section{User Story 4 - 保存 ETF 研究結果}
\textbf{Story ID:} \texttt{US-MKT-04}
\textbf{用戶故事：} 作為完成初步研究的用戶，我希望將 ETF 保存到 \texttt{Watchlist} 繼續跟蹤，或加入已有 \texttt{My Portfolio} 做進一步組合分析，並可通過統一 \texttt{Screener} 繼續篩選。
\textbf{Video:} \href{https://jjpvro70sief.jp.larksuite.com/wiki/WE34wt8oyiO3NOkBA3zjFrxppge}{Market 驗收視頻頁面}
\subsection{Walkthrough}
\begin{enumerate}
\item 在 \texttt{ETF Detail} 點擊 star icon，將 ETF 加入 \texttt{Watchlist}。
\item 打開 \texttt{Screener} 的 \texttt{Watchlist} 標簽，確認該 ETF 出現；隨後移除並確認列表同步更新。
\item 回到 \texttt{ETF Detail}，點擊 plus button，打開 \texttt{Add to existing portfolio} sheet。
\item 展示 selected ETF、latest price、lots 和 estimated market value。說明 1 lot 按 100 shares 計算。
\item 展示 duplicate holding 狀態：已含該 ETF 的 portfolio 顯示重複提示且不可選擇；再選擇另一個 eligible portfolio 保存。
\item 重新打開目標 \texttt{My Portfolio}，確認 ETF 已作為 holding record 保存。強調該流程不下單、不執行交易。
\item 返回 \texttt{Screener}，依次展示 \texttt{Recommended}、\texttt{All}、\texttt{Watchlist} 三個資料範圍，並操作一次 search、filter 和 sort。
\end{enumerate}
\screenshotsTripleCompact{assets/screenshots/mkt-us4-01-watchlist-saved.png}{MKT-US4-01}{\texttt{Watchlist ETFs} 保存結果，顯示已選 \texttt{Watchlist} 標簽、保存數量及 ETF symbol/\allowbreak{}name。}{assets/screenshots/mkt-us4-02-add-existing-portfolio.png}{MKT-US4-02}{\texttt{Add to existing} sheet，包含 selected ETF、lots、estimated value、duplicate/\allowbreak{}disabled portfolio 和 eligible portfolio。}{assets/screenshots/mkt-us4-03-screener-controls.png}{MKT-US4-03}{統一 \texttt{Screener} 的 \texttt{All /\allowbreak{} Recommended /\allowbreak{} Watchlist} scopes、search、filter、sort metric 和 return period controls。}
\screenshotDescsThree{\textbf{MKT-US4-01} \textemdash{} \texttt{Watchlist ETFs} 保存結果，顯示已選 \texttt{Watchlist} 標簽、保存數量及 ETF symbol/\allowbreak{}name。}{\textbf{MKT-US4-02} \textemdash{} \texttt{Add to existing} sheet，包含 selected ETF、lots、estimated value、duplicate/\allowbreak{}disabled portfolio 和 eligible portfolio。}{\textbf{MKT-US4-03} \textemdash{} 統一 \texttt{Screener} 的 \texttt{All /\allowbreak{} Recommended /\allowbreak{} Watchlist} scopes、search、filter、sort metric 和 return period controls。}
\subsection{Acceptance Criteria}
\begingroup
\small
\begin{longtable}{L{0.18\linewidth}L{0.30\linewidth}L{0.18\linewidth}L{0.12\linewidth}L{0.10\linewidth}}
\toprule
\textbf{ID} & \textbf{Requirement} & \textbf{Reference} & \textbf{Ownership} & \textbf{Priority} \\
\midrule
MKT-US4-AC01 & 點擊 \texttt{ETF Detail} 的 star icon 後，成功狀態與 \texttt{Watchlist} 中的實際保存狀態一致。 & US-MKT-04 & Integration & Must \\
MKT-US4-AC02 & 已保存 ETF 出現在 \texttt{Screener \textgreater{} Watchlist}；移除後從該列表消失。 & US-MKT-04 & Integration & Must \\
MKT-US4-AC03 & \texttt{Watchlist} 用於保存研究對象，不會創建 portfolio holding 或交易訂單。 & US-MKT-04 & Integration & Must \\
MKT-US4-AC04 & plus button 為當前 ETF 打開 \texttt{Add to existing portfolio} sheet。 & US-MKT-04 & Integration & Must \\
MKT-US4-AC05 & sheet 展示 selected ETF、latest price、lots 和 estimated market value，並按 1 lot = 100 shares 計算。 & US-MKT-04 & Integration & Must \\
MKT-US4-AC06 & 已包含該 ETF 的 \texttt{My Portfolio} 顯示 duplicate holding 信息且不可再次選擇。 & US-MKT-04 & Integration & Must \\
MKT-US4-AC07 & 用戶可以選擇另一個 eligible \texttt{My Portfolio} 並完成保存。 & US-MKT-04 & Integration & Must \\
MKT-US4-AC08 & 保存完成後，重新打開目標 \texttt{My Portfolio} 可看到對應 ETF holding；未選擇並確認前不會寫入。 & US-MKT-04 & Integration & Must \\
MKT-US4-AC09 & \texttt{Recommended}、\texttt{All}、\texttt{Watchlist} 共用同一 \texttt{Screener} 交互，但分別限定為精選池、完整可用範圍和用戶已保存範圍。 & US-MKT-04 & Integration & Must \\
MKT-US4-AC10 & search 支持 symbol、name 或 issuer；filter 支持當前實現提供的 issuer、asset class、sector、currency 或 region。 & US-MKT-04 & Integration & Must \\
MKT-US4-AC11 & sort 支持當前實現提供的 change、volume、turnover、price、expense ratio、beta、drawdown、Sharpe ratio、name 或 symbol。 & US-MKT-04 & Integration & Must \\
MKT-US4-AC12 & search、filter 和 sort 不會意外切換當前 \texttt{Recommended}、\texttt{All} 或 \texttt{Watchlist} 資料範圍。 & US-MKT-04 & Integration & Must \\
MKT-US4-AC13 & \texttt{Add to My Portfolio} 只保存研究型 holding record，頁面不聲稱已經下單或執行交易。 & US-MKT-04 & Integration & Must \\
\bottomrule
\end{longtable}
\endgroup
